\section{Medidas Robustas}

Tomemos o exemplo de uma base que contenha salários de funcionários de uma empresa.
A média pode ser pouco representativa para descrever esses dados, pois, normalmente, há
poucos funcionários que ganham altos salários, como executivos, enquanto a grande maioria
dos funcionários tem uma remuneração consideravelmente mais baixa.
Por consequência, a variância e o desvio-padrão também podem não ser apropriados para
medir a variabilidade, pois os valores extremos podem indicar uma dispersão enquanto que
os dados, na realidade, concentram-se num determinado intervalo.

A mediana e o IQR, usados para medir centralidade e dispersão respectivamente, sofrem
menos influência da magnitude do dado do que a média e variância, por exemplo.
Isso se deve ao fato de que seus valores são computados utilizando somente parte
dos dados, e não por uma agregação completa deles.
Logo, na suspeita da presença de valores extremos e de forte assimetria, essas medidas
podem ser mais apropriadas para realizar a descrição da base.

Há também outras medidas muito utilizas para descrever dados de maneira robusta.
A primeira a citar é a \textbf{média aparada}, que consiste no cálculo da média
considerando apenas parte dos dados menos extremos.
Para isso, decidi-se por um parâmetro $\alpha \in [0, 1]$ e calcula-se a média apenas
das estatísticas entre aquelas que acumulam $\alpha$ e $1-\alpha$ de probabilidade.
Isso garante que a média não será influenciada por valores extremos, sendo, pois, mais
representativa.

Uma segunda medida, dessa vez de dispersão, é o \textbf{coeficiente de variação}.
Ele é calculado apenas tomando-se o desvio-padrão e o divindido pela média, produzindo,
por conseguinte, uma grandeza adimensional.
Essa propriedade permite avaliar a variabilidade de um atributo de forma padronizada, ou
seja, independente da escala dele, tornando possível, também, comparar comparar a
variabilidade de dois atributos de magnitude ou escala distintos.

Uma segunda medida robusta de dispersão é o \textbf{desvio absoluto mediano}, análogo ao
desvio absoluto médio.
Para o computar, basta tomar a mediana dos desvios das estatísticas em relação a mediana
do atributo.
Assim, valores extremos pouco interferem no valor final.
